% !TEX root = ../my-thesis.tex

\chapter{Fazit und Ausblick}
\label{sec:conclusion}

\section{Fazit}
\label{sec:conclusion:fazit}

\begin{sloppypar}
Diese Arbeit hat fünf Workflow-Management-Systeme anhand zweier
Workflow-Muster und derselben Rechenlogik verglichen, zunächst
unter kontrollierten Bedingungen auf einem einzelnen Rechner und
anschließend auf dem HPC-System MOGON. Untersucht wurden die
Modellierung, das Laufzeitverhalten anhand von Makespan, Overhead
und Koordinationsgrößen sowie der Betrieb unter Slurm.
\end{sloppypar}

Beide Muster ließen sich in allen fünf Systemen umsetzen. Die
wesentlichen Unterschiede liegen daher nicht in der Abbildbarkeit,
sondern in der Art der Workflow-Beschreibung, im
Koordinationsverhalten und im Betriebsmodell. Beim Overhead
bildeten StreamFlow, Nextflow und Merlin auf beiden Ebenen eine
Gruppe mit niedrigen Werten, AiiDA lag darüber und Pegasus in der
lokalen Untersuchung mit großem Abstand an der Spitze. Mit zunehmender Rechenlast verliert der Overhead gegenüber der
eigentlichen Rechenzeit an Bedeutung.

Auf dem HPC-System zeigte sich ein starker Einfluss der
Ausführungsumgebung auf die gemessenen Laufzeiten. Die vom Cluster verursachten Zeitanteile können die Unterschiede zwischen den Systemen überlagern, und die Systeme binden den Ressourcenverwalter
unterschiedlich ein, von einzelnen Jobs je Schritt bis zur
gemeinsamen Worker-Allokation. Für den Systemvergleich in dieser Arbeit wurde deshalb innerhalb einer bestehenden Allokation gemessen.

Für die Eignungsfrage folgt daraus, dass sie sich nicht an der
Abbildbarkeit der Muster festmachen lässt, denn diese war in allen
fünf Systemen gegeben. Unterschiede zeigten sich dagegen im
Laufzeitverhalten und im Betrieb. Mit abnehmender Bedeutung des
Overheads können für die Wahl daher andere, aus der Ausrichtung der
Systeme hervorgehende Eigenschaften stärker ins Gewicht fallen.
Dazu zählen beispielsweise die Unterstützung hybrider
Ausführungsumgebungen bei StreamFlow, die Ausrichtung von Merlin
auf große Ensembles von Rechenaufgaben oder die dauerhafte
Nachvollziehbarkeit bei AiiDA. Eine allgemeine Rangfolge der
Systeme ergibt sich daraus nicht. Welches System geeignet
erscheint, hängt damit auch davon ab, welche dieser Eigenschaften
für den jeweiligen Anwendungsfall im Vordergrund stehen.

\section{Ausblick}
\label{sec:conclusion:ausblick}

An den Grenzen der Arbeit setzen mögliche weiterführende
Arbeiten an. Der Vergleich ließe sich zunächst auf komplexere
Workflow-Strukturen ausdehnen und mit Rechenlasten wiederholen, die
realen wissenschaftlichen Anwendungen näherkommen, etwa mit großen
Datenmengen und intensiven Dateizugriffen. So ließe sich prüfen, ob
die hier beobachteten Unterschiede zwischen den Systemen auch unter
solchen Bedingungen fortbestehen.

Ein zweiter Ansatzpunkt ist die Skalierung. Die Untersuchung
umfasste höchstens vier parallele Verarbeitungsinstanzen. Mit
deutlich mehr Aufgaben ließe sich prüfen, wie sich die Systeme bei
großen Aufgabenzahlen verhalten, insbesondere Merlin, dessen Entwurf
genau darauf zielt. Messungen auf weiteren HPC-Systemen würden
zudem zeigen, welche der auf MOGON beobachteten Betriebs- und
Laufzeiteigenschaften an die konkrete Ausführungsumgebung gebunden
sind.

Offen bleibt schließlich der HPC-Vergleich für Pegasus. Lässt sich
Pegasus mit der benötigten HTCondor-Umgebung auf einem Cluster
betreiben, kann die quantitative Untersuchung auf alle fünf Systeme
ausgeweitet werden.